\section{Problemas encontrados y comentarios}

\subsection{Incidencias reales}

\begin{table}[H]
  \centering
  \begin{tabularx}{\textwidth}{>{\raggedright\arraybackslash\bfseries}p{0.23\textwidth}XX}
    \toprule
    Incidencia & Causa & Solución aplicada \\
    \midrule
    Acceso denegado al socket de Docker & Aunque el usuario ya pertenecía al
    grupo \texttt{docker}, la sesión activa todavía no había cargado dicha
    pertenencia. & Se aplicó el grupo en una subsesión con \texttt{sg docker};
    un nuevo inicio de sesión produce el mismo efecto. \\
    \addlinespace
    Instrucciones no específicas para Fedora & La guía proporcionaba ejemplos
    para Debian, Ubuntu, Windows y macOS. & Se conservaron los objetivos de la
    práctica y se utilizaron los paquetes y el gestor de servicios propios de
    Fedora. \\
    \bottomrule
  \end{tabularx}
  \caption{Problemas observados y soluciones empleadas.}
\end{table}

\subsection{Comentarios finales}

El uso de contenedores permitió desplegar PostgreSQL sin incorporar archivos
del servidor al sistema operativo anfitrión. El volumen independiente conserva
los datos aunque el contenedor se detenga o se sustituya, mientras que el mapeo
limitado a \texttt{localhost} reduce la superficie de exposición durante el
trabajo local.

La validación se efectuó en tres niveles: ejecución de un contenedor mínimo,
consulta directa con \texttt{psql} y conexión mediante DBeaver. Esta secuencia
permitió distinguir problemas del motor Docker, del servidor PostgreSQL y del
cliente gráfico. Al concluir, la credencial temporal en texto plano se eliminó;
DBeaver conservó únicamente su representación dentro del almacén seguro de la
aplicación.

\begin{nota}{Resultado}
Docker quedó habilitado y activo; PostgreSQL 18.6 permaneció en ejecución bajo
el nombre \texttt{postgres-practica01}; y DBeaver Community 26.2.0 se conectó
correctamente a la base \texttt{postgres} mediante JDBC.
\end{nota}
